\section*{C4: Logica Hoare}

\subsection{Exercițiul 1}

Regula consecinței

\[
	\frac{\models A \rightarrow A' \quad \{A'\} c \{B'\} \quad \models B' \rightarrow B}{\{A\} c \{B\}}
\]

Este o \textbf{pseudo-regulă}, fiindcă nu este decidabil dacă o aplicare a ei
este corectă.

\textbf{Justificare.}
Pentru a verifica o aplicare trebuie să decidem validitatea formulelor
$A\rightarrow A'$ și $B'\rightarrow B$. Validitatea formulelor din logica de
ordinul I este indecidabilă.

Mai precis, fie $\varphi$ o aserțiune arbitrară. Avem
\[
	\frac{\models \mathit{true}\rightarrow\varphi
		\quad \{\varphi\}\ \instr{skip}\ \{\varphi\}
		\quad \models\varphi\rightarrow\mathit{true}}
	{\{\mathit{true}\}\ \instr{skip}\ \{\mathit{true}\}}.
\]
Premisa din mijloc este o instanță a axiomei pentru $\instr{skip}$, iar ultima
este mereu validă. Prima este validă exact atunci când
$\models\varphi$. Deci, dacă am putea decide dacă această aplicare a regulii
este corectă, am putea decide validitatea oricărei formule $\varphi$, ceea ce
este imposibil.

Regula rămâne însă \textbf{corectă}: dacă cele trei premise sunt adevărate,
atunci și concluzia este adevărată. „Pseudo” se referă doar la faptul că
aplicabilitatea ei nu poate fi verificată algoritmic în general.

\subsection{Exercițiul 2}

Pentru orice $A,B,c,I$, demonstrăm că
\[
	\models^I\{A\}c\{B\}
	\iff
	A^I\subseteq ws^I(c,B),
\]
unde
\[
	A^I=\{\sigma\in\Sigma\mid \sigma\models^I A\}
\]
și
\[
	ws^I(c,B)=
	\{\sigma\in\Sigma\mid
	\forall\sigma'\,((\sigma,\sigma')\in\llbracket c\rrbracket
	\Rightarrow \sigma'\models^I B)\}.
\]

\textbf{$(\Rightarrow)$}
Presupunem $\models^I\{A\}c\{B\}$ și fie $\sigma\in A^I$.
Atunci $\sigma\models^I A$. Din semantica enunțului Hoare, pentru orice
$\sigma'$ cu $(\sigma,\sigma')\in\llbracket c\rrbracket$, avem
$\sigma'\models^I B$. Deci $\sigma\in ws^I(c,B)$.

\textbf{$(\Leftarrow)$}
Presupunem $A^I\subseteq ws^I(c,B)$. Fie
$(\sigma,\sigma')\in\llbracket c\rrbracket$ și presupunem
$\sigma\models^I A$. Atunci $\sigma\in A^I$, deci
$\sigma\in ws^I(c,B)$. Prin definiția lui $ws^I$, rezultă
$\sigma'\models^I B$. Așadar $\models^I\{A\}c\{B\}$.


\subsection{Exercițiul 3}
Exercițiul din pagina 23 din cursul 4.

Presupunem că limbajul permite disjuncții infinitare și că
$w:=\mathbf{while}\ b\ \mathbf{do}\ c$ se termină din orice stare. Definim
\[
	P_0:=\neg b\land B,\qquad
	P_{k+1}:=b\land wp(c,P_k),
\]
și
\[
	wp(w,B):=\bigvee_{k\in\mathbb N}P_k.
\]

Pentru incluziunea $ws^I(w,B)\subseteq wp(w,B)^I$, fie
$\sigma\in ws^I(w,B)$. Cum $w$ se termină, există o execuție finită
\[
	\sigma_0=\sigma,\ \sigma_1,\ldots,\sigma_n
\]
astfel încât, pentru $i<n$,
\[
	\sigma_i\models^I b,\qquad
	(\sigma_i,\sigma_{i+1})\in\llbracket c\rrbracket,
\]
iar $\sigma_n\models^I\neg b\land B$.

\textbf{Ce demonstrăm.}
Prin inducție după $j$, pentru orice $0\leq j\leq n$,
\[
	\sigma_{n-j}\models^I P_j.
\]

\textbf{Cazul de bază: $j=0$.}
Avem $\sigma_n\models^I\neg b\land B=P_0$, deoarece bucla s-a oprit în
$\sigma_n$, iar $\sigma\in ws^I(w,B)$ garantează postcondiția $B$.

\textbf{Pasul inductiv.}
Presupunem, pentru $j<n$, că
$\sigma_{n-j}\models^I P_j$. Notăm $i:=n-j-1$. Din descrierea execuției,
\[
	\sigma_i\models^I b,\qquad
	(\sigma_i,\sigma_{i+1})\in\llbracket c\rrbracket,
	\qquad \sigma_{i+1}\models^I P_j.
\]
Semantica lui $c$ este deterministă; deci orice stare finală obținută
executând $c$ din $\sigma_i$ este $\sigma_{i+1}$ și satisface $P_j$.
Prin definiția celei mai slabe precondiții,
\[
	\sigma_i\models^I wp(c,P_j).
\]
Împreună cu $\sigma_i\models^I b$, rezultă
\[
	\sigma_i\models^I b\land wp(c,P_j)=P_{j+1}.
\]
Cum $i=n-(j+1)$, pasul inductiv este demonstrat.

Pentru $j=n$ obținem
\[
	\sigma=\sigma_0\models^I P_n,
\]
deci $\sigma\models^I\bigvee_{k\in\mathbb N}P_k=wp(w,B)$.
